\chapter{涨落理论与复习路线}

\section{正则系综中的能量涨落}

正则系综中
\[
  U=\avg{E}=-\frac{\partial}{\partial\beta}\ln Z.
\]
能量涨落为
\[
  \avg{(\Delta E)^2}
  =\avg{E^2}-\avg{E}^2
  =\frac{\partial^2}{\partial\beta^2}\ln Z.
\]
利用热容定义可得
\[
  \avg{(\Delta E)^2}=\kb T^2 C_V.
\]
这说明热容本质上是能量涨落的响应。

\section{巨正则系综中的粒子数涨落}

巨正则系综中
\[
  \avg{N}
  =\frac{1}{\beta}\pdc{\ln\Xi}{\mu}{T,V}.
\]
粒子数涨落为
\[
  \avg{(\Delta N)^2}
  =\frac{1}{\beta^2}\frac{\partial^2 \ln\Xi}{\partial\mu^2}
  =\kb T\pdc{\avg{N}}{\mu}{T,V}.
\]
对普通宏观系统，
\[
  \frac{\sqrt{\avg{(\Delta N)^2}}}{N}\sim N^{-1/2}.
\]

\section{涨落-响应关系}

许多响应函数都可看作涨落：
\[
  C_V=\frac{\avg{(\Delta E)^2}}{\kb T^2},
\]
\[
  \kappa_T
  \propto \frac{\avg{(\Delta N)^2}}{N\kb T}
  \quad \text{或与体积涨落相关}.
\]
接近临界点时，相关长度增大，涨落增强，响应函数可能发散。这是本科相变与高等统计物理临界现象之间的桥梁。

\section{本科热统到高统的复习路线}

建议把这门基础课压缩成四条线：
\begin{enumerate}
  \item 热力学势线：$U,F,G,H,\Omega$ 及其自然变量；
  \item 计数线：$\micro,W,Z,\Xi$ 与熵、自由能；
  \item 分布线：Boltzmann, Bose, Fermi；
  \item 涨落线：热容、压缩率、粒子数涨落和临界行为。
\end{enumerate}

如果目标是服务高等统计物理，最优先掌握：
\[
  S=\kb\ln\micro,\qquad
  p_i=\frac{e^{-\beta E_i}}{Z},\qquad
  F=-\kb T\ln Z,
\]
\[
  \Xi=\sum_N e^{\beta\mu N}Z_N,\qquad
  \Phi=-\kb T\ln\Xi,\qquad
  \avg{(\Delta E)^2}=\kb T^2 C_V.
\]

\begin{keybox}
\textbf{复习判断标准：}看到一个体系，能说出它的宏观约束、该选哪个系综、配分函数怎么写、自由能怎么取、涨落对应哪个响应函数，就已经具备进入高统的基础。
\end{keybox}
